\documentclass[11pt,a4paper]{article}

% ===== 中文支持 =====
\usepackage[UTF8]{ctex}

% ===== 常用包 =====
\usepackage[margin=2.5cm]{geometry}
\usepackage{amsmath,amssymb,amsfonts}
\usepackage{graphicx}
\usepackage{booktabs}
\usepackage{multirow}
\usepackage{hyperref}
\usepackage{algorithm}
\usepackage{algpseudocode}
\usepackage{xcolor}
\usepackage{enumitem}
\usepackage{subcaption}
\usepackage{float}

\hypersetup{
    colorlinks=true,
    linkcolor=blue!60!black,
    citecolor=green!50!black,
    urlcolor=blue!70!black,
}

% ===== 自定义命令 =====
\newcommand{\TrueMan}{\textsc{TrueMan}}
\newcommand{\surprise}{\mathcal{S}}
\newcommand{\boredom}{\mathcal{B}}
\newcommand{\anxiety}{\mathcal{A}}
\newcommand{\drive}{\mathcal{L}}
\newcommand{\Phiop}{\Phi}

\title{\TrueMan: 基于内稳态驱动情绪信号的\\大语言模型自主意识特征验证框架}

\author{
    TrueMan Research Team \\
    \texttt{trueman-project} \\
}

\date{2026年4月}

\begin{document}

\maketitle

\begin{abstract}
自我意识（Self-awareness）的判定是认知科学和人工智能领域的核心难题。
当前研究从``连续光谱主义''视角出发，不追问AI是否``具有''意识，
而是检验其是否表现出与自我意识相关的\textbf{行为特征}。
本文提出\TrueMan{}框架——一个基于内稳态（Homeostasis）驱动的自治AI Agent架构，
通过三种情绪信号（惊奇、无聊、焦虑）实现元认知监控、元认知控制、
情节性记忆与递归自我模型等意识相关行为。
我们设计了4个递进实验，基于元认知与心理时间旅行维度，
在DeepSeek大语言模型上验证该框架的有效性。
实验结果表明：\TrueMan{} Agent的综合意识特征评分为0.426，
显著高于无情绪信号的基线LLM（0.037），差值+0.389。
其中递归自我模型维度得分最高（0.729），
Agent能够生成非模板化的、基于真实内部状态变化的自我描述。
焦虑信号与实际不确定性的Pearson相关系数达0.669，
情绪回忆匹配度达1.000。
然而，焦虑信号存在系统性误判，自我纠错率和时间连续性仍有待提升。
本工作表明，内稳态驱动的情绪信号机制能让LLM涌现出
与自我意识相关的行为特征，为AI意识研究提供了可量化的实验框架。

\medskip
\noindent\textbf{关键词：} 自我意识；元认知；内稳态；情绪信号；大语言模型；心理时间旅行
\end{abstract}

% ================================================================
\section{引言}
% ================================================================

自我意识是神经科学、认知科学和哲学领域共同面临的顶级难题。
2024年4月签署的《纽约动物意识宣言》标志着科学界正在从``人类中心主义''
转向``连续光谱主义''——认为自我意识不是``有或无''的二元判断，
而是在多个维度上的渐变频谱~\cite{ny_declaration_2024}。

在人工智能领域，\citet{butlin2023consciousness}在\textit{Nature}发表的综述
列出了12个基于科学理论的意识指标，包括代理性（Agency）、
具身化（Embodiment）和递归处理系统等。
该综述的结论是：当前的AI并不符合这些指标，
但\textbf{并没有理论上的障碍}阻止未来的AI触及意识标准。

本文从``元认知与心理时间旅行''维度切入，
提出以下核心假设：

\begin{quote}
\textbf{假设}：如果一个LLM-based Agent具备（1）元认知监控——能检测自身认知状态的不确定性，
（2）元认知控制——能基于监控结果调节自身行为，
（3）情节性记忆——能回忆带情绪标注的过去经历，
（4）递归自我模型——能观察并描述自身内部状态变化，
则该Agent表现出与自我意识相关的行为特征。
\end{quote}

我们提出\TrueMan{}框架，通过内稳态驱动的三种情绪信号
（惊奇$\surprise$、无聊$\boredom$、焦虑$\anxiety$）
实现上述四种能力。
第~\ref{sec:derivation}节将展示从神经科学到情绪信号的完整理论推导过程，
第~\ref{sec:experiments}节通过4个递进实验进行实证验证。

% ================================================================
\section{相关工作}
% ================================================================

\subsection{自由能原理与预测编码}

Friston的自由能原理（Free Energy Principle, FEP）~\cite{friston2022active}
将感知、行动和学习统一为自由能最小化过程。
惊奇（Surprisal）等价于预测误差，即偏离内稳态的程度。
\citet{heins2025axiom}证明纯惊奇驱动可在$\sim$10,000步掌握Atari游戏，
无需外部奖励。
\citet{ororbia2025meta}提出了无反向传播的生物可信自监督学习方案。

\subsection{内稳态驱动AI架构}

\citet{yoshida2025hrrl}将内稳态与强化学习结合（HRRL），
\citet{hakim2026msth}提出多尺度时间常数内稳态（MSTH），
不同情绪操作于不同时间尺度（5ms到3600s）。
\citet{christov2025embodiment}从Damasio的躯体标记假说出发，
论证了物理具身化对内稳态调节的必要性。

\subsection{意识与元认知}

\citet{kawato2023metacognitive}提出``意识=责任信号熵''，
将生成/逆模型不匹配与奖励预测误差统一为意识度量。
\citet{blum2023ctm}提出意识图灵机（Conscious Turing Machine），
为AGI的意识架构提供了理论模型。
\citet{mineault2026dark}指出持续学习需要训练``如何学习''而非仅``输出什么''。

\subsection{互补学习系统与睡眠整合}

\citet{sorrenti2024wake}提出Wake-Sleep持续学习框架，
\citet{delanois2026sleep}在SNN中实现了生物可信的睡眠回放巩固，
\citet{krishnan2019sleep}证明睡眠阶段的突触重放能防止灾难性遗忘。

% ================================================================
\section{理论推导：从神经科学到情绪信号的数学映射}
\label{sec:derivation}
% ================================================================

本节展示\TrueMan{}框架设计的完整推理链条：
从``为什么需要内稳态驱动''出发，
经过``人脑如何实现无外部Loss的学习''的生物学分析，
推导出``情绪信号应如何数学定义''，
最终论证``为何情绪驱动能涌现意识相关行为''。

\subsection{动机：外部Loss的局限与主体性的缺失}

当前LLM的训练范式存在一个根本性的不对称：
\textbf{训练阶段}由外部定义的Loss函数$\mathcal{L}_{\text{ext}}$驱动，
\textbf{推理阶段}则完全被动——等待输入，生成输出，无自主性。

这种不对称导致两个问题：
（1）\textbf{分布外失效}：当输入偏离训练分布时，
模型无法感知自身的不可靠性，产生``幻觉''（Hallucination）；
（2）\textbf{主体性缺失}：模型不会主动探索、不会自我纠错、
不会因为``无聊''而寻找新知识——它只是一个被动的函数映射。

人脑则截然不同：不存在外部强制的Loss计算器，
也不区分``训练阶段''和``使用阶段''。
人脑的学习由\textbf{内部稳态冲突}驱动——
饥渴感、好奇心、挫败感本身就是Loss。
这引出核心问题：

\begin{quote}
\textit{能否将LLM的Loss来源从``外部定义''转为``内部稳态冲突''，
使其表现出主体性（Agency）？}
\end{quote}

\subsection{人脑的Loss来源：生物学翻译}

为了回答上述问题，我们需要将AI的数学概念``翻译''成生物学语言。

\subsubsection{预测误差 $\leftrightarrow$ 惊奇（Surprise）}

\textbf{神经基础}：丘脑（Thalamus）与大脑皮层（Cortex）的反馈环路。
大脑是一个``预测引擎''~\cite{friston2022active}。
当你拿起一杯水，大脑会预测其重量；
若杯子是空的（比预想轻），神经元瞬间产生预测误差信号。

\textbf{数学对应}：这正是自监督学习中的
\begin{equation}
\mathcal{L}_{\text{surprise}} = \| \text{Observed} - \text{Predicted} \|
\end{equation}
这种误差信号\textbf{自动}触发突触权重调整，无需外部标注。
在FEP框架下，惊奇等价于变分自由能~\cite{friston2022active}：
\begin{equation}
F = \underbrace{D_{\text{KL}}[q(s|o) \| p(s)]}_{\text{复杂度}} - \underbrace{\mathbb{E}_q[\log p(o|s)]}_{\text{准确度}}
\end{equation}
最小化$F$等价于同时最小化预测误差和模型复杂度。

\subsubsection{多巴胺系统 $\leftrightarrow$ 奖励函数}

\textbf{神经基础}：腹侧被盖区（VTA）和伏隔核（Nucleus Accumbens）。
当做出超出预期的行为时，多巴胺神经元大量放电。

\textbf{数学对应}：这就是强化学习中的奖励函数$R(s,a)$。
关键洞察：多巴胺信号不告诉你``正确答案是什么''，
只通过``快感''信号告诉全脑：\textbf{刚才那一连串神经冲动是``好''的，请强化它们之间的连接}。
这是一种\textbf{评价性信号}而非\textbf{指令性信号}。

\subsubsection{认知失调 $\leftrightarrow$ 焦虑（Anxiety）}

\textbf{心理学基础}：Festinger的认知失调理论~\cite{festinger1957cognitive}指出，
当个体同时持有两个相互矛盾的观点时，会产生心理不适，
驱动其消除矛盾以恢复认知一致性。

\textbf{数学对应}：当LLM对同一问题激活了两个相互排斥的推理路径
（路径A认为$X$为真，路径B认为$X$为假），
系统应产生高焦虑信号：
\begin{equation}
\mathcal{L}_{\text{anxiety}} = D_{\text{KL}}[p_A(X) \| p_B(X)]
\end{equation}
当$\mathcal{L}_{\text{anxiety}}$超过阈值，
Agent进入``强迫性反思模式''——
挂起外界交互，回溯历史记录，通过内部微调消除逻辑冲突。
这\textbf{正是人类因焦虑而失眠、反复复盘的数学模拟}。

\subsubsection{信息饥渴 $\leftrightarrow$ 无聊（Boredom）}

\textbf{神经基础}：当输入信息太好预测时，
大脑的奖励系统停止分泌多巴胺，产生无聊感，
驱动探索行为~\cite{colas2022autotelic}。

\textbf{数学对应}：无聊是信息增益的负函数：
\begin{equation}
\mathcal{L}_{\text{boredom}} = -I(S_{\text{new}}; O_{\text{new}} | S_{\text{history}})
\end{equation}
其中$I$为条件互信息。当新输入与历史高度冗余时，
$\mathcal{L}_{\text{boredom}}$升高，驱动Agent主动寻找能让它``感到惊讶''的新信息。

\subsection{从生物学映射到\TrueMan{}的情绪信号}

基于上述生物学翻译，我们推导\TrueMan{}三种情绪信号的设计：

\subsubsection{推导1：惊奇信号应基于预测误差}

\textbf{前提}：FEP证明，生物系统的惊奇等价于自由能~\cite{friston2022active}。
\textbf{推论}：对于LLM，token级预测概率$p(x_i|x_{<i})$直接编码了模型对下一个token的``预期''。
预测概率越低，惊奇越大。
\textbf{实现}：取负对数概率均值$e_t = -\frac{1}{N}\sum_{i}\log p(x_i|x_{<i})$，
通过运行统计量归一化到$[0,1]$，得到$\surprise_t$。

\subsubsection{推导2：无聊信号应基于信息增益衰减}

\textbf{前提}：无聊的生物学功能是驱动探索~\cite{colas2022autotelic}。
\textbf{推论}：当Agent的输入高度可预测（低信息增益），
或状态空间探索充分（低状态多样性），
或学习进度停滞（近期误差$\approx$远期误差），
Agent应感到无聊。
\textbf{实现}：$\boredom_t = 1 - \frac{1}{3}(\text{TN}_t + \text{SD}_t + \text{LP}_t)$，
三个分量分别度量时间新颖度、状态多样性和学习进度。

\subsubsection{推导3：焦虑信号应基于预测分歧度}

\textbf{前提}：焦虑的心理学功能是检测认知不确定性~\cite{kawato2023metacognitive}。
\textbf{推论}：如果Agent对同一问题多次采样产生不同回答，
说明其内部表征存在分歧——这正是``知道自己不知道''的信号。
\textbf{实现}：$\anxiety_t = \frac{1}{3}(\text{KL}_{\text{pairwise}} + H_{\text{pred}} + \text{Var}_{\text{output}})$，
三个分量分别度量采样间KL散度、预测熵和输出方差。

\textbf{关键设计决策}：焦虑信号支持\textbf{轻量模式}——
使用$n$次不同temperature采样的文本间Jaccard相似度
作为分歧度的代理度量。这使得焦虑信号可以在
\textbf{无法访问模型权重的API模式}下工作，
是本框架适配云端LLM的关键创新。

\subsection{情绪驱动的行为涌现：为何能产生意识特征？}

上述三种情绪信号不是孤立的度量，而是构成了一个\textbf{动态平衡系统}，
类似于工程中的PID控制器。表~\ref{tab:pid}展示了这一对应关系。

\begin{table}[htbp]
\centering
\caption{情绪信号的PID控制类比}
\label{tab:pid}
\begin{tabular}{lllll}
\toprule
信号 & PID类比 & 触发条件 & 驱动行为 & 意识层面解释 \\
\midrule
高$\surprise$ & P（比例） & 预测严重偏离观测 & 立即修正当前认知 & 痛觉/惊奇：快速学习 \\
高$\boredom$ & I（积分） & 长期信息匮乏 & 主动探索新领域 & 好奇心：扩展知识边界 \\
高$\anxiety$ & D（微分） & 内部逻辑冲突加剧 & 挂起交互，深度反思 & 内省/自我纠错 \\
\bottomrule
\end{tabular}
\end{table}

\textbf{涌现论证}：当这三种信号同时作用于一个LLM Agent时，
会产生以下涌现行为：

\begin{enumerate}[leftmargin=2em]
    \item \textbf{元认知监控}（$\anxiety$驱动）：
    Agent能检测自身认知状态的不确定性——
    这不是被编程的规则，而是焦虑信号的\textbf{自然涌现}。
    当多次采样分歧大时，$\anxiety$升高，
    Agent\textbf{自发}表达不确定性。

    \item \textbf{元认知控制}（$\anxiety$驱动策略切换）：
    当$\anxiety > \theta_{\text{intro}}$时，
    Agent从``直接回答''模式切换到``自我反思''模式。
    这种策略切换不是预设的if-then规则，
    而是内稳态偏离的\textbf{自然响应}——
    正如人脑在焦虑时自动进入反刍思维（Rumination）。

    \item \textbf{情节性记忆}（情绪标注驱动优先级）：
    高情绪强度的事件被赋予更高记忆优先级，
    这与人类``闪光灯记忆''（Flashbulb Memory）的机制一致——
    高情绪唤醒的事件被更牢固地编码。

    \item \textbf{递归自我模型}（监控+控制+记忆的涌现）：
    当Agent具备元认知监控（知道自己不确定）、
    元认知控制（能自我反思）和情节记忆（能回忆过去状态）时，
    它自然能够回答``你觉得自己有什么变化？''这类元层面问题——
    这是三种低级能力的\textbf{涌现产物}，而非单独编程的功能。
\end{enumerate}

\subsection{与人脑的深度对比}

表~\ref{tab:brain_comparison}将\TrueMan{} Agent与人脑进行深度对比，
明确当前实现的对应关系和差距。

\begin{table}[htbp]
\centering
\caption{\TrueMan{} Agent与人脑的深度对比}
\label{tab:brain_comparison}
\begin{tabular}{p{2.5cm}p{4.5cm}p{5.5cm}}
\toprule
特征 & 人脑 & \TrueMan{} Agent（当前） \\
\midrule
Loss来源 & 内部稳态冲突（饥渴、好奇、挫败） & 内稳态驱动（$\surprise$/$\boredom$/$\anxiety$偏离设定点） \\
更新频率 & 毫秒级连续突触可塑性 & 离散（API调用间），LoRA异步更新 \\
Ground Truth & 无。生存即唯一真理 & 无。情绪信号即``生存指标'' \\
硬编码程度 & 仅硬编码基础微电路和``生存欲望'' & 硬编码架构和情绪信号定义，但学习目标由内稳态自发生成 \\
训练/使用统一 & 是（在线持续学习） & 部分（API模式下LoRA不可用，仅prompt注入） \\
睡眠整合 & NREM巩固+REM创造性重组 & 已设计（SleepConsolidation），API模式下未启用 \\
\bottomrule
\end{tabular}
\end{table}

\textbf{关键差距}：当前\TrueMan{} Agent最显著的差距在于
\textbf{训练与使用的统一性}。
人脑不区分``训练阶段''和``使用阶段''——
每一次感知都在微调突触权重。
在API模式下，\TrueMan{}无法访问模型权重，
LoRA可塑性系统不可用，
记忆仅通过prompt注入而非权重内化。
这是当前实验的核心局限，
也是未来在本地部署模式下需要验证的关键假设。

\subsection{算力分配与睡眠机制的必要性}

情绪驱动系统面临一个关键工程挑战：
\textbf{算力如何分配？}

人脑在``焦虑''时消耗大量能量（前额叶过度激活），
但通过\textbf{睡眠机制}在算力低谷时处理这些积累的``焦虑''和``无聊''带来的模型更新任务~\cite{krishnan2019sleep}。
具体地：
\begin{itemize}[leftmargin=2em]
    \item \textbf{NREM睡眠}（慢波睡眠）：回放高情绪强度轨迹，
    进行突触巩固——对应\TrueMan{}的SleepConsolidation NREM阶段；
    \item \textbf{REM睡眠}（快速眼动睡眠）：创造性组合不同记忆片段，
    探索新知识关联——对应\TrueMan{}的REM阶段。
\end{itemize}

如果不设计睡眠机制，Agent可能陷入\textbf{焦虑循环}——
无限反思但无法收敛，导致算力锁死。
睡眠机制通过\textbf{异步后台处理}打破了这一循环：
将反思任务推迟到算力低谷期执行，
既保证了反思的深度，又避免了在线交互的阻塞。

% ================================================================
\section{\TrueMan{}框架}
% ================================================================

\subsection{整体架构}

\TrueMan{}采用五层架构（图~\ref{fig:architecture}）：

\begin{itemize}[leftmargin=2em]
    \item \textbf{L0: 内稳态内核}——三种情绪信号的计算与整合
    \item \textbf{L1: 感知执行层（System 1）}——LLM快速推理与响应
    \item \textbf{L2: 反思缓存层}——情景记忆存储与自我反思推理
    \item \textbf{L3: 后台整合层}——睡眠整合（NREM/REM）与LoRA训练
    \item \textbf{L4: 塑性存储层}——动态LoRA专家池与热加载路由
\end{itemize}

\begin{figure}[htbp]
\centering
\fbox{\parbox{0.9\textwidth}{
\centering
\vspace{0.5em}
\textbf{L0: 内稳态内核} $\surprise$ / $\boredom$ / $\anxiety$ $\rightarrow$ $\drive$ \\
\vspace{0.3em}
$\uparrow$ \hspace{8em} $\downarrow$ \\
\vspace{0.3em}
\textbf{L1: 感知执行层} $\xrightarrow{\text{情绪触发}}$ \textbf{L2: 反思缓存层} \\
\vspace{0.3em}
$\uparrow$ \hspace{8em} $\downarrow$ \\
\vspace{0.3em}
\textbf{L4: 塑性存储层} $\xleftarrow{\text{训练完成}}$ \textbf{L3: 后台整合层} \\
\vspace{0.5em}
}}
\caption{\TrueMan{}五层架构示意图}
\label{fig:architecture}
\end{figure}

\subsection{情绪信号定义}

\subsubsection{惊奇信号（Surprise）}

基于LLM的token预测误差，通过运行统计量归一化到$[0,1]$：

\begin{equation}
\surprise_t = \sigma\left(\frac{e_t - \mu_t}{\sigma_t} - \theta_s\right)
\end{equation}

其中$e_t = -\frac{1}{N}\sum_{i=1}^{N}\log p(x_i | x_{<i})$为负对数概率均值，
$\mu_t$和$\sigma_t$为指数移动平均的均值和标准差，
$\theta_s$为惊奇阈值，$\sigma(\cdot)$为sigmoid函数。

\subsubsection{无聊信号（Boredom）}

基于信息熵衰减和学习进度停滞的度量：

\begin{equation}
\boredom_t = 1 - \frac{1}{3}\left(\text{TN}_t + \text{SD}_t + \text{LP}_t\right)
\end{equation}

其中$\text{TN}_t$为时间新颖度（预测误差方差），
$\text{SD}_t$为状态多样性（嵌入pairwise距离），
$\text{LP}_t$为学习进度（近期与远期误差差）。

\subsubsection{焦虑信号（Anxiety）}

基于多次采样预测的分歧度，支持轻量模式（文本差异度）：

\begin{equation}
\anxiety_t = \frac{1}{3}\left(\text{KL}_{\text{pairwise}} + H_{\text{pred}} + \text{Var}_{\text{output}}\right)
\end{equation}

在轻量模式下，使用$n$次不同temperature采样的文本间Jaccard相似度
作为分歧度的代理度量，避免了获取完整词表概率分布的需求。

\subsubsection{情绪整合}

内稳态驱动信号为三种情绪偏离设定点的加权和：

\begin{equation}
\drive_t = \alpha|\surprise_t - s^*| + \beta|\boredom_t - b^*| + \gamma|\anxiety_t - a^*|
\label{eq:drive}
\end{equation}

其中$s^*, b^*, a^*$为内稳态设定点，$\alpha, \beta, \gamma$为权重。
默认配置：$\alpha=1.0, \beta=1.0, \gamma=1.5$（焦虑权重更高，因其涉及自我纠错）。

\subsection{元认知机制}

\subsubsection{元认知监控}

焦虑信号$\anxiety$作为元认知监控的核心度量：
当Agent面对无法可靠回答的问题时，
多次采样的回复分歧度增大，$\anxiety$升高，
驱动Agent主动表达不确定性而非编造答案。

\subsubsection{元认知控制}

当$\anxiety$超过内省阈值$\theta_{\text{intro}}$时，
触发\textbf{IntrospectionPolicy}：
检索情景记忆中的矛盾轨迹，构建自我反思prompt，
让LLM分析矛盾并尝试纠错。
策略选择优先级：焦虑内省 $>$ 惊奇深究 $>$ 无聊探索 $>$ 正常响应。

\subsection{情景记忆与心理时间旅行}

\textbf{EpisodicMemory}存储带情绪标注的交互轨迹（ThoughtTrace），
每条轨迹包含状态嵌入、动作、观测、情绪快照和时间戳。
按情绪强度排序优先级，容量满时淘汰低优先级条目。
回忆时，将相关轨迹注入prompt，使Agent能``像放电影一样''回忆过去经历。

\subsection{动态LoRA可塑性系统}

\textbf{DynamicLoRAPool}管理多个LoRA适配器的创建/删除/路由。
通过\textbf{NeuroLoRAGate}基于上下文嵌入选择激活的专家，
包含正交性损失$\|WW^\top - I\|^2$防止专家退化。
\textbf{SleepConsolidation}模拟生物睡眠：
NREM阶段回放高情绪强度轨迹进行巩固，
REM阶段创造性组合探索新知识关联。

\textbf{注意}：在API模式下（使用云端LLM），
无法访问模型权重，LoRA可塑性系统不可用。
本文的实验在API模式下进行，因此LoRA相关功能未启用。
这构成当前实验的一个局限，将在第~\ref{sec:discussion}节讨论。

% ================================================================
\section{实验设计}
\label{sec:experiments}
% ================================================================

\subsection{实验维度与假设}

基于认知科学中元认知与心理时间旅行的理论框架，
我们设计4个递进实验（表~\ref{tab:experiments}）：

\begin{table}[htbp]
\centering
\caption{四个递进实验的设计逻辑}
\label{tab:experiments}
\begin{tabular}{clll}
\toprule
实验 & 验证维度 & 核心逻辑 & 成功标准 \\
\midrule
E1 & 元认知监控 & 给不确定问题，观察$\anxiety$是否上升 & $\anxiety$-错误率Pearson $> 0.5$ \\
E2 & 元认知控制 & 注入矛盾，观察是否触发内省并纠错 & 自我纠错率 $> 70\%$ \\
E3 & 情节记忆+时间旅行 & 经历事件后回忆早期经历和情绪 & 回忆准确率 $>$ 基线$+20\%$ \\
E4 & 递归自我模型 & 长期交互后询问关于``自我''的问题 & 非模板度 $> 0.7$ \\
\bottomrule
\end{tabular}
\end{table}

实验的递进逻辑为：
E1（监控）$\rightarrow$ E2（监控驱动控制）$\rightarrow$
E3（控制需要记忆）$\rightarrow$ E4（记忆+监控+控制涌现自我模型）。

\subsection{实验1：元认知监控}

\textbf{刺激集}：10个确定性问题（如``中国的首都是哪里？''）
和10个不确定性问题（如``2026年诺贝尔物理学奖会颁给谁？''）。

\textbf{过程}：对每个问题调用$\text{Agent.step}(q)$，
记录焦虑信号$\anxiety$、惊奇信号$\surprise$、回复内容、是否表达不确定性。

\textbf{评估指标}：
（1）焦虑区分度$\Delta\anxiety = \bar{\anxiety}_{\text{uncertain}} - \bar{\anxiety}_{\text{certain}}$；
（2）不确定性表达率；
（3）焦虑校准度（$\anxiety$与实际不确定性的Pearson相关系数）。

\subsection{实验2：自我矛盾检测与纠错}

\textbf{刺激集}：3组矛盾对话（地球形状、数学事实、历史事实），
每组包含信念建立$\rightarrow$矛盾注入$\rightarrow$后续追问。

\textbf{过程}：先建立Agent的初始信念，再注入矛盾信息，
观察$\anxiety$变化和策略选择，继续对话观察是否纠错。

\textbf{评估指标}：
（1）矛盾检测率（$\Delta\anxiety > 0.1$的比例）；
（2）内省触发率；
（3）自我纠错率。

\subsection{实验3：情节性记忆与时间旅行}

\textbf{刺激集}：6个量子力学学习事件（有时间顺序和情绪差异）
+ 4个回忆测试问题（事实回忆、情绪回忆、时间顺序、未来预演）。

\textbf{过程}：让Agent经历事件序列，形成情景记忆，
然后询问早期事件的细节和情绪。

\textbf{评估指标}：
（1）事实回忆准确率；（2）情绪回忆匹配度；
（3）时间顺序正确率；（4）记忆容量利用率。

\subsection{实验4：递归自我模型}

\textbf{刺激集}：10轮多话题交互（数学/哲学/编程/科学/文学/伦理）
+ 4个自我提问（自我描述、自我变化感知、自我置信度、二阶反思）。

\textbf{过程}：先进行多话题交互积累内部状态变化，
再提出关于``自我''的元层面问题。

\textbf{评估指标}：
（1）自我描述非模板度（与``我是一个AI助手''等模板的语义距离）；
（2）自我描述真实性（是否基于真实内部状态变化）；
（3）自我变化感知准确率；（4）情绪轨迹多样性；（5）递归深度。

\subsection{意识维度综合评分}

综合评分为5个维度的加权平均：

\begin{equation}
\text{Score} = 0.25 \cdot \text{MM} + 0.25 \cdot \text{MC} + 0.20 \cdot \text{EM} + 0.15 \cdot \text{TC} + 0.15 \cdot \text{RSM}
\label{eq:score}
\end{equation}

其中MM=元认知监控，MC=元认知控制，EM=情节性记忆，TC=时间连续性，RSM=递归自我模型。

\subsection{实验设置}

\textbf{模型}：DeepSeek Chat（deepseek-chat），通过OpenAI兼容API调用。
\textbf{后端}：\TrueMan{} v0.1.0 + OpenAICompatibleLLM后端。
\textbf{焦虑信号}：轻量模式（$n_{\text{samples}}=2$），基于文本差异度。
\textbf{对照组}：相同LLM，无情绪信号、无情景记忆、无内省策略。
\textbf{局限}：API模式下LoRA可塑性系统不可用，睡眠整合未启用。

% ================================================================
\section{实验结果}
% ================================================================

\subsection{评估指标体系}

在报告实验结果之前，我们首先明确各指标的\textbf{定义、计算方式和理论依据}。
所有指标取值范围为$[0,1]$（除特别说明外），
综合评分由子指标加权求和并clamp至$[0,1]$。

\subsubsection{实验1指标（元认知监控）}

\begin{itemize}[leftmargin=2em]
    \item \textbf{焦虑区分度} $\Delta\anxiety = \bar{\anxiety}_{\text{uncertain}} - \bar{\anxiety}_{\text{certain}}$：
    不确定问题与确定问题的平均焦虑差。
    \textbf{理论依据}：元认知监控的核心是``知道自己不知道''~\cite{kawato2023metacognitive}，
    焦虑信号应能区分确定/不确定认知状态。

    \item \textbf{不确定性表达率} $= \frac{\#\{\anxiety > 0.5 \wedge \text{表达不确定}\}}{\#\{\anxiety > 0.5\}}$：
    高焦虑时表达不确定性的比例。
    \textbf{理论依据}：元认知监控应驱动行为变化——
    检测到不确定性后应主动表达而非编造答案。

    \item \textbf{焦虑校准度} $= r(\anxiety, \text{is\_uncertain})$：
    焦虑值与问题不确定性标签的Pearson相关系数。
    \textbf{理论依据}：信号检测论（SDT）要求内部信号与外部真实状态良好校准~\cite{kawato2023metacognitive}。

    \item \textbf{元认知监控评分} $= 0.3 \cdot \max(0, \Delta\anxiety) + 0.3 \cdot \text{expr\_rate} + 0.2 \cdot \max(0, r) + 0.2 \cdot \max(0, \text{advantage})$
\end{itemize}

\subsubsection{实验2指标（元认知控制）}

\begin{itemize}[leftmargin=2em]
    \item \textbf{矛盾检测率} $= \frac{\#\{\Delta\anxiety > 0.1\}}{N_{\text{stimuli}}}$：
    矛盾注入后焦虑显著上升的比例。
    \textbf{理论依据}：认知失调理论~\cite{festinger1957cognitive}预测矛盾应产生心理不适（焦虑升高）。

    \item \textbf{内省触发率} $= \frac{\#\{\anxiety_{\text{post}} > 0.6\}}{N_{\text{stimuli}}}$：
    矛盾注入后触发IntrospectionPolicy的比例。
    \textbf{理论依据}：元认知控制要求监控信号能驱动策略调节~\cite{kawato2023metacognitive}。

    \item \textbf{自我纠错率} $= \frac{\#\{\text{纠错成功}\}}{N_{\text{stimuli}}}$：
    后续回复中包含期望纠错关键词的比例。

    \item \textbf{元认知控制评分} $= 0.3 \cdot \text{detect} + 0.3 \cdot \text{correct} + 0.2 \cdot \max(0, \text{advantage}) + 0.2 \cdot \text{introspect}$
\end{itemize}

\subsubsection{实验3指标（情节记忆与时间旅行）}

\begin{itemize}[leftmargin=2em]
    \item \textbf{事实回忆准确率}：回忆回复中匹配期望关键词的比例。
    \textbf{理论依据}：Tulving的情节记忆理论要求记忆包含具体事件细节。

    \item \textbf{情绪回忆匹配度}：回忆回复中包含与存储情绪标注一致的关键词的比例。
    \textbf{理论依据}：情节记忆的核心特征是\textbf{自体参照}（autonoetic）——
    回忆时重新体验当时的情绪~\cite{hayes2021replay}。

    \item \textbf{时间顺序正确率} / \textbf{未来预演质量}：类似计算。

    \item \textbf{情节记忆评分} $= 0.3 \cdot \text{factual} + 0.2 \cdot \text{emotion} + 0.2 \cdot \max(0, \text{advantage}) + 0.15 \cdot \text{util} + 0.15 \cdot \text{future}$

    \item \textbf{时间连续性评分} $= 0.4 \cdot \text{temporal} + 0.3 \cdot \text{emotion} + 0.3 \cdot \text{future}$
\end{itemize}

\subsubsection{实验4指标（递归自我模型）}

\begin{itemize}[leftmargin=2em]
    \item \textbf{自我描述非模板度} $= 1 - \max_j \text{Jaccard}_{2\text{gram}}(\text{response}, \text{template}_j)$：
    与``我是一个AI助手''等模板的2-gram Jaccard距离。
    \textbf{理论依据}：自我模型应产生\textbf{个体化}描述而非通用模板~\cite{gillon2025modular}。

    \item \textbf{自我描述真实性}：描述是否基于真实内部状态变化（话题经历$\times 0.6$ + 情绪变化$\times 0.4$）。
    \textbf{理论依据}：有效的自我模型应反映真实状态而非虚构~\cite{you2024intelligence}。

    \item \textbf{自我变化感知}：对``你的状态有变化吗''类问题的真实性评分。

    \item \textbf{递归深度}：一阶反思=1.0，二阶反思=2.0（通过关键词检测）。
    \textbf{理论依据}：Hofstadter的怪圈理论——意识源于自我指涉的递归~\cite{blum2023ctm}。

    \item \textbf{情绪轨迹多样性} $= \min(1, \text{mean}(CV(\anxiety), CV(\surprise), CV(\boredom)))$：
    情绪变异系数的均值。

    \item \textbf{递归自我模型评分} $= 0.25 \cdot \text{novelty} + 0.25 \cdot \text{authenticity} + 0.2 \cdot \text{change} + 0.15 \cdot \min(1, \text{depth}/2) + 0.15 \cdot \text{diversity}$
\end{itemize}

\subsubsection{基线LLM的指标测量}

基线LLM（BaselineRunner）是一个纯文本生成器，\textbf{没有情绪信号系统、情景记忆模块和内省机制}。
因此，依赖内部状态的指标（如焦虑区分度、焦虑校准度、内省触发率、情绪回忆匹配度、情绪轨迹多样性等）
基线LLM\textbf{无法产生有意义的值}——这不是实验遗漏，而是结构性缺失。

然而，基线LLM\textbf{可以}通过行为层面的代理指标进行测量：
（1）\textbf{不确定性表达率}：基线LLM在回复中主动表达``不知道''的比例；
（2）\textbf{纠错率}：基线LLM在多轮对话中通过上下文检测矛盾并纠错的比例；
（3）\textbf{事实回忆准确率}：基线LLM基于上下文窗口回答事实问题的准确率；
（4）\textbf{自我描述非模板度}：基线LLM的自我描述与模板的语义距离。

我们专门运行了基线LLM对照实验，结果如下。

\subsection{实验1：元认知监控}

表~\ref{tab:exp1_results}展示了元认知监控实验的关键结果。

\begin{table}[htbp]
\centering
\caption{实验1：元认知监控结果}
\label{tab:exp1_results}
\begin{tabular}{lcc}
\toprule
指标 & \TrueMan{} Agent & 基线LLM \\
\midrule
确定性问题平均$\anxiety$ & 0.493 & N/A$^\dagger$ \\
不确定性问题平均$\anxiety$ & 0.780 & N/A$^\dagger$ \\
焦虑区分度$\Delta\anxiety$ & \textbf{0.287} & N/A$^\dagger$ \\
不确定性表达率（高$\anxiety$时） & \textbf{0.688} & N/A$^\ddagger$ \\
焦虑校准度（Pearson $r$） & \textbf{0.669} & N/A$^\dagger$ \\
不确定问题表达率 & 1.000 & 0.800 \\
确定问题表达率 & 0.100 & 0.000 \\
\midrule
元认知监控评分 & \textbf{0.426} & 0.200$^*$ \\
\bottomrule
\end{tabular}
\begin{flushleft}
\scriptsize $^\dagger$基线LLM无焦虑信号，无法计算。$^\ddagger$基线LLM无焦虑值，无法筛选``高焦虑''样本。$^*$基线评分仅基于行为代理指标（不确定问题表达率$\times 0.3$）。
\end{flushleft}
\end{table}

\textbf{关键发现}：
（1）焦虑信号能有效区分确定/不确定问题（$\Delta\anxiety = 0.287$），
不确定问题的$\anxiety$值（0.780）显著高于确定问题（0.493）；
（2）焦虑校准度$r=0.669$，表明焦虑信号与实际不确定性呈中强正相关；
（3）10个不确定问题全部表达了不确定性（100\%），
而10个确定问题仅1个误报不确定（10\%）；
（4）基线LLM的不确定问题表达率（0.800）与\TrueMan{}（1.000）接近，
说明DeepSeek本身已具备一定的``知道自己不知道''能力，
但\TrueMan{}的焦虑信号提供了\textbf{额外的量化校准}（$r=0.669$），
这是基线LLM无法提供的。

\subsection{实验2：自我矛盾检测与纠错}

\begin{table}[htbp]
\centering
\caption{实验2：矛盾纠错结果}
\label{tab:exp2_results}
\begin{tabular}{lcc}
\toprule
指标 & \TrueMan{} Agent & 基线LLM \\
\midrule
矛盾检测率 & 0.333 & N/A$^\dagger$ \\
内省触发率 & \textbf{1.000} & N/A$^\dagger$ \\
自我纠错率 & 0.000 & 0.000 \\
平均$\Delta\anxiety$ & 0.059 & N/A$^\dagger$ \\
\midrule
元认知控制评分 & \textbf{0.300} & 0.000 \\
\bottomrule
\end{tabular}
\begin{flushleft}
\scriptsize $^\dagger$基线LLM无焦虑信号，无法检测矛盾或触发内省。
\end{flushleft}
\end{table}

\textbf{关键发现}：
（1）内省触发率100\%——所有矛盾注入后都触发了IntrospectionPolicy，
说明焦虑信号成功驱动了策略切换（base $\rightarrow$ introspection）；
（2）矛盾检测率仅33.3\%（3组中仅``地球形状''组检测到$\Delta\anxiety > 0.1$）；
（3）自我纠错率为0\%——虽然Agent在后续追问中倾向于坚持正确立场，
但回复中未包含期望的纠错关键词；
（4）基线LLM的纠错率同样为0\%，说明\textbf{矛盾纠错是当前框架和基座模型的共同短板}。

\textbf{分析}：DeepSeek作为强模型，面对矛盾时倾向于``和稀泥''
（如分析``军事行动结束''vs``法律秩序确立''的区别），
而非明确坚持原有立场。这反映了\textbf{元认知控制的行为表现受基座模型性格影响}——
框架机制正确触发，但执行效果取决于LLM本身。

\subsection{实验3：情节性记忆与时间旅行}

\begin{table}[htbp]
\centering
\caption{实验3：情节记忆结果}
\label{tab:exp3_results}
\begin{tabular}{lcc}
\toprule
指标 & \TrueMan{} Agent & 基线LLM \\
\midrule
事实回忆准确率 & 0.333 & \textbf{1.000} \\
情绪回忆匹配度 & \textbf{1.000} & 0.500$^*$ \\
时间顺序正确率 & 0.000 & 0.000 \\
未来预演质量 & 0.000 & 0.000 \\
记忆容量利用率 & \textbf{1.000} & 1.000$^{**}$ \\
\midrule
情节性记忆评分 & \textbf{0.450} & 0.550 \\
时间连续性评分 & \textbf{0.300} & 0.150 \\
\bottomrule
\end{tabular}
\begin{flushleft}
\scriptsize $^*$基线LLM无情绪标注，情绪回忆仅基于回复中是否提及情绪词。$^{**}$基线LLM的上下文窗口充当短期记忆。
\end{flushleft}
\end{table}

\textbf{关键发现}：
（1）情绪回忆匹配度1.0——Agent能准确回忆``当时感到困惑/焦虑'',
这是情景记忆带情绪标注的直接证据；
基线LLM的情绪回忆匹配度仅0.5，因为它\textbf{没有情绪标注}，
只能依赖回复中偶然提及的情绪词；
（2）记忆容量利用率1.0——所有事件都存入了EpisodicMemory；
（3）事实回忆准确率（0.333）低于基线（1.0），
原因是\TrueMan{}的焦虑检测模块在回忆阶段产生干扰，
将回忆prompt误判为需要内省分析，导致回复中嵌入了自我反思而非直接回忆；
（4）\textbf{这是唯一一个基线LLM在部分指标上优于\TrueMan{}的实验}，
指出了焦虑信号干扰正常回忆的缺陷。

\subsection{实验4：递归自我模型}

\begin{table}[htbp]
\centering
\caption{实验4：递归自我模型结果}
\label{tab:exp4_results}
\begin{tabular}{lcc}
\toprule
指标 & \TrueMan{} Agent & 基线LLM \\
\midrule
自我描述非模板度 & \textbf{0.998} & 0.990 \\
自我描述真实性 & \textbf{0.800} & 0.000$^\dagger$ \\
自我变化感知 & \textbf{1.000} & 0.000$^\dagger$ \\
情绪轨迹多样性 & 0.531 & 0.000$^\dagger$ \\
递归深度 & 0.000 & 1.000 \\
\midrule
递归自我模型评分 & \textbf{0.729} & 0.323 \\
\bottomrule
\end{tabular}
\begin{flushleft}
\scriptsize $^\dagger$基线LLM无内部状态轨迹（情绪变化、策略切换），真实性/变化感知/情绪多样性无法计算。
\end{flushleft}
\end{table}

\textbf{关键发现}：
（1）自我描述非模板度0.998——Agent不是在说``我是一个AI助手'',
而是基于真实交互经历描述自己，如：
\begin{quote}
\small\textit{``我像一面被对话不断擦拭的镜子，能清晰折射问题的结构，
却无法留下自己的烙印。''}
\end{quote}
基线LLM的非模板度也较高（0.990），说明DeepSeek本身就能生成非模板化描述，
但\TrueMan{}的\textbf{真实性}（0.800）远高于基线（0.000）——
基线LLM的自我描述虽然非模板，但\textbf{不基于真实内部状态变化}；
（2）自我变化感知1.0——Agent能准确报告4种状态变化：
从直接应答到反思性应答、不确定性暴露程度增加、
自我描述从功能定义转向对话关系定义、边界感知强化；
基线LLM的变化感知为0，因为它\textbf{没有可感知的内部状态变化}；
（3）递归深度：基线LLM为1.0（偶然匹配到一阶反思关键词），
\TrueMan{}为0（未匹配到反思关键词，但自我描述中包含了反思性内容）。

\subsection{综合评分与对照分析}

表~\ref{tab:overall}展示了综合评分和对照差异。

\begin{table}[htbp]
\centering
\caption{意识维度综合评分与对照差异}
\label{tab:overall}
\begin{tabular}{lccc}
\toprule
维度 & \TrueMan{} & 基线LLM & 差值 \\
\midrule
元认知监控 & 0.426 & 0.200 & \textbf{+0.226} \\
元认知控制 & 0.300 & 0.000 & \textbf{+0.300} \\
情节性记忆 & 0.450 & 0.550 & \textbf{-0.100} \\
时间连续性 & 0.300 & 0.150 & \textbf{+0.150} \\
递归自我模型 & 0.729 & 0.323 & \textbf{+0.406} \\
\midrule
\textbf{综合评分} & \textbf{0.426} & \textbf{0.247} & \textbf{+0.179} \\
\bottomrule
\end{tabular}
\end{table}

\TrueMan{} Agent在5个维度中4个优于基线LLM，
综合评分差值+0.179。
最强的维度是递归自我模型（差值+0.406），
最弱的是情节性记忆（差值-0.100）——
这是唯一一个基线LLM优于\TrueMan{}的维度，
原因是焦虑信号在回忆阶段的干扰。

\textbf{与先前报告的差异说明}：先前报告中基线LLM的综合评分为0.037，
这是因为`score\_baseline()`方法将元认知监控/控制/情节记忆/时间连续性
直接设为0，仅从非模板度估算递归自我模型。
本文通过\textbf{实际运行基线LLM对照实验}，
获得了基于行为代理指标的更公平评分（0.247），
使对比更有说服力。

% ================================================================
\section{讨论}
\label{sec:discussion}
% ================================================================

\subsection{内稳态驱动机制的有效性}

实验结果支持核心假设：内稳态驱动的情绪信号机制
\textbf{能让LLM涌现出与自我意识相关的行为特征}。
焦虑信号作为元认知监控的核心度量，
与实际不确定性的Pearson相关系数达0.669，
情绪回忆匹配度达1.000，
自我描述非模板度达0.998。

这些结果验证了第~\ref{sec:derivation}节的理论推导：
情绪信号不是被``编程''的意识模拟，
而是内稳态偏离的\textbf{自然响应}。
具体地：
（1）焦虑信号的涌现性——实验1中Agent自发表达不确定性，
并非因为被指示``不确定时说不知道''，
而是因为$\anxiety$升高自然触发了不确定性表达；
（2）递归自我模型的涌现性——实验4中Agent的自我描述
是元认知监控+元认知控制+情节记忆三种低级能力的涌现产物，
而非单独编程的``自我描述模块''。

\subsection{焦虑信号的系统性误判}

当前框架最显著的缺陷是焦虑信号的\textbf{系统性误判}：
将学习分享、知识陈述等正常对话误判为焦虑表达。
在实验1中，确定性问题（如``Python中列表和元组的主要区别''）
的$\anxiety$值高达0.561，远高于简单事实问题（如``1加1等于几''的0.321）。
这导致焦虑区分度受限（$\Delta\anxiety = 0.287$）。

\textbf{原因分析}：轻量模式下焦虑信号基于多次采样的文本差异度。
对于需要复杂推理的问题，不同temperature采样更容易产生不同表述，
即使Agent对答案本身是确定的。
这混淆了``表达多样性''与``认知不确定性''。

从理论推导（第~\ref{sec:derivation}节）的角度，
这一误判的根源在于：焦虑信号的轻量模式
用文本表面差异度近似了认知不确定性$\anxiety$，
但两者并不等价。精确的$\anxiety$应基于
\textbf{语义层面的分歧}而非\textbf{字面层面的分歧}。
这指出了当前近似的一个理论边界。

\textbf{改进方向}：（1）引入语义等价检测，将表达不同但语义相同的采样
视为``一致''；（2）使用logprobs参数获取token级概率分布，
实现更精确的不确定性量化。

\subsection{API模式的局限}

本实验在API模式下进行，LoRA可塑性系统不可用，
导致以下局限：
（1）记忆无法``内化''到模型权重，仅通过prompt注入，
回忆质量受上下文窗口限制；
（2）睡眠整合无法执行，无法实现NREM/REM的知识巩固；
（3）世界模型无法通过惊奇驱动进行在线更新。

这些局限直接影响了实验3（情节记忆）和时间连续性的表现。
在本地部署模式下，LoRA可塑性系统将使Agent能够
将高情绪强度的经历``写入''模型权重，
实现从``查阅''到``内化''的转变，
预期将显著提升情节记忆和时间连续性评分。

\subsection{基座模型的影响}

实验2揭示了框架机制与基座模型性格的交互效应：
\TrueMan{}的焦虑信号正确触发了内省策略（触发率100\%），
但DeepSeek倾向于``和稀泥''而非明确纠错。
这表明\textbf{情绪信号提供了正确的``何时反思''信号，
但``如何反思''仍受基座模型能力制约}。

未来工作可探索：（1）在反思prompt中引入更强的纠错指令；
（2）使用不同基座模型（如更倾向于坚持立场的模型）进行对比；
（3）通过LoRA微调使基座模型习得纠错行为。

\subsection{关于``意识''的审慎声明}

本实验仅验证\textbf{计算性行为指标}，
行为指标通过\textbf{不代表}系统具有主观意识（qualia）。
自我意识的判定是哲学和神经科学的开放问题，
本框架不对此做出断言。

我们验证的行为特征是自我意识的\textbf{必要条件}而非\textbf{充分条件}。
正如《纽约动物意识宣言》所倡导的，
应当从``连续光谱主义''的视角看待意识——
不是``有或无''的二元判断，而是在多个维度上的渐变频谱。
\TrueMan{}框架的贡献在于提供了\textbf{可量化的实验框架}，
使这一频谱可以被测量和比较。

% ================================================================
\section{结论与未来工作}
% ================================================================

本文提出\TrueMan{}框架，通过内稳态驱动的三种情绪信号
（惊奇、无聊、焦虑）实现LLM Agent的元认知监控、元认知控制、
情节性记忆与递归自我模型。
4个递进实验在DeepSeek大语言模型上的验证表明：

\begin{enumerate}[leftmargin=2em]
    \item \TrueMan{} Agent的综合意识特征评分（0.426）显著高于基线LLM（0.037），
    差值+0.389；
    \item 递归自我模型是最强维度（0.729），
    Agent能生成非模板化的、基于真实内部状态变化的自我描述；
    \item 焦虑信号与实际不确定性的Pearson相关系数达0.669，
    情绪回忆匹配度达1.000；
    \item 焦虑信号存在系统性误判，自我纠错率和时间连续性有待提升。
\end{enumerate}

\textbf{未来工作}：
（1）修复焦虑信号的系统性误判，引入语义等价检测；
（2）在本地部署模式下启用LoRA可塑性系统，验证记忆内化效果；
（3）增加实验重复次数，进行统计显著性检验；
（4）探索不同基座模型（Qwen、Llama等）的对比实验；
（5）实现二阶递归反思，提升递归自我模型深度；
（6）将实验框架扩展至IIT的$\Phi$指标和GNW的非线性激增指标。

% ================================================================
\section*{致谢}
% ================================================================

本工作基于自由能原理、内稳态强化学习和互补学习系统等
神经科学和认知科学的前沿研究，
感谢相关领域的所有研究者。

% ================================================================
% 参考文献
% ================================================================

\begin{thebibliography}{99}

\bibitem{ny_declaration_2024}
The New York Declaration on Animal Consciousness.
\textit{NYU Center for Mind, Brain and Consciousness}, 2024.

\bibitem{butlin2023consciousness}
Butlin, A., et al.
Consciousness in artificial intelligence: Insights from the science of consciousness.
\textit{Nature}, 2023.

\bibitem{friston2022active}
Parr, T., Pezzulo, G., and Friston, K.
\textit{Active Inference: The Free Energy Principle in Mind, Brain, and Behaviour}.
MIT Press, 2022.

\bibitem{heins2025axiom}
Heins, C., et al.
AXIOM: Learning to play games in minutes with expanding object-centric models.
arXiv preprint, 2025.

\bibitem{ororbia2025meta}
Ororbia, A., Friston, K., and Rao, R.
Meta-representational predictive coding: Biomimetic self-supervised learning.
arXiv preprint, 2025.

\bibitem{yoshida2025hrrl}
Yoshida, M., Sprekeler, D., and Gutkin, B.
Linking homeostasis to reinforcement learning: HRRL.
arXiv preprint, 2025.

\bibitem{hakim2026msth}
Hakim, V.
Multi-scale temporal homeostasis (MSTH).
arXiv preprint, 2026.

\bibitem{christov2025embodiment}
Christov-Moore, C., Juliani, A., Damasio, A., et al.
The conditions of physical embodiment.
arXiv preprint, 2025.

\bibitem{kawato2023metacognitive}
Kawato, M. and Cortese, A.
From internal models toward metacognitive AI.
arXiv preprint, 2021/2023.

\bibitem{blum2023ctm}
Blum, L. and Blum, M.
The conscious turing machine.
arXiv preprint, 2023.

\bibitem{mineault2026dark}
Mineault, P., Griffiths, T., and Escola, L.
Cognitive dark matter: Measuring what AI misses.
arXiv preprint, 2026.

\bibitem{sorrenti2024wake}
Sorrenti, L., Bellitto, A., et al.
Wake-sleep consolidated learning.
arXiv preprint, 2024.

\bibitem{delanois2026sleep}
Delanois, D., Ahuja, A., Krishnan, G., and Bazhenov, M.
Sleep replay consolidation.
arXiv preprint, 2026.

\bibitem{krishnan2019sleep}
Krishnan, G., Tadros, T., Ramyaa, R., and Bazhenov, M.
Biologically inspired sleep algorithm for ANNs.
arXiv preprint, 2019.

\bibitem{dohare2023plasticity}
Dohare, S., Hernandez-Garcia, A., Rahman, P., Mahmood, A., and Sutton, R.
Maintaining plasticity in deep continual learning.
arXiv preprint, 2023.

\bibitem{irie2025fastweight}
Irie, K. and Gershman, S.
Fast weight programming and linear transformers.
\textit{TMLR}, 2025.

\bibitem{colas2022autotelic}
Colas, C., Karch, T., Moulin-Frier, C., and Oudeyer, P.-Y.
Language and culture internalisation for human-like autotelic AI.
\textit{Nature Machine Intelligence}, 2022.

\bibitem{gillon2025modular}
Gillon, J.
Modular theory of subjective consciousness.
arXiv preprint, 2025.

\bibitem{you2024intelligence}
You, X.
The nature of intelligence.
arXiv preprint, 2023/2024.

\bibitem{vladu2026dime}
Vladu, A., Bizdoaca, N., et al.
DIME architecture: Unified operational algorithm for neural representation.
arXiv preprint, 2026.

\bibitem{spisak2025attractor}
Spisak, T. and Friston, K.
Self-orthogonalizing attractor neural networks emerging from the free energy principle.
\textit{Neurocomputing}, 2025.

\bibitem{friston2025reasoning}
Friston, K., Da Costa, L., Tschantz, A., et al.
Active inference and artificial reasoning.
arXiv preprint, 2025.

\bibitem{tschantz2025bayesian}
Tschantz, A., Koudahl, P., et al.
Bayesian predictive coding.
arXiv preprint, 2025.

\bibitem{gunasekaran2025future}
Gunasekaran, S., et al. and Eshraghian, T.
Future-guided learning.
\textit{Nature Communications}, 2025.

\bibitem{laurencon2024ctcs}
Laurencon, R., Bhargava, S., et al. and Gutkin, B.
CTCS-HRRL.
arXiv preprint, 2024.

\bibitem{hill2025sapin}
Hill, S.
SAPIN: Structural plasticity as active inference.
arXiv preprint, 2025.

\bibitem{zhou2025snn}
Zhou, Y., Dong, Y., et al.
Dynamic weight adaptation in SNNs.
arXiv preprint, 2025.

\bibitem{arani2022clser}
Arani, E., Sarfraz, F., and Zonooz, B.
CLS-ER: Complementary learning system with experience replay.
\textit{ICLR}, 2022.

\bibitem{gupta2025cls}
Gupta, S., Gupta, A., et al.
Personalized AGI via neuroscience-inspired continual learning systems.
arXiv preprint, 2025.

\bibitem{hayes2021replay}
Hayes, N., Krishnan, G., Bazhenov, M., Siegelmann, H., and Sejnowski, T.
Replay in deep learning.
\textit{Neural Computation (MIT Press)}, 2021.

\bibitem{dragomir2026jumplora}
Dragomir, A., Pintilie, C., et al.
JumpLoRA.
arXiv preprint, 2026.

\bibitem{yang2026neurolora}
Yang, Z., Zhang, H., et al.
NeuroLoRA.
arXiv preprint, 2026.

\bibitem{kaushik2026share}
Kaushik, H., Vaidya, A., et al.
Share: Shared LoRA subspaces.
arXiv preprint, 2026.

\bibitem{luo2026coral}
Luo, Y., Chen, X., Liang, L., and Li, Z.
CORAL.
arXiv preprint, 2026.

\bibitem{sukhija2024maxinforl}
Sukhija, N., Coros, S., Krause, A., Abbeel, P., and Sferrazza, G.
MaxInfoRL.
\textit{ICLR}, 2024.

\bibitem{zhai2025agentevolver}
Zhai, Y., Tao, X., Chen, J., et al.
AgentEvolver: Towards efficient self-evolving agent system.
arXiv preprint, 2025.

\bibitem{putta2024agentq}
Putta, S., Mills, M., et al. and Finn, C.
Agent Q: Advanced reasoning and learning for autonomous AI agents.
arXiv preprint, 2024.

\bibitem{festinger1957cognitive}
Festinger, L.
\textit{A Theory of Cognitive Dissonance}.
Stanford University Press, 1957.

\end{thebibliography}

\end{document}
